% Context from chapters/sparse-regularization.tex; for mathematical reference, not compilation.
\section{Sparsity Priors}

                                                                      
\subsection{Ideal Sparsity Prior}

Chapter \ref{chap-approximation} shows that an orthonormal basis $\Bb = \{ \psi_m \}_m$ adapted to an image class $\Th$ can approximate each $f\in\Th$ using relatively few atoms.

The $\lzero$ prior measures representation complexity by counting the nonzero coefficients:
\eq{
	\Jz(f) \eqdef \#\enscond{m}{\dotp{f}{\psi_m} \neq 0}
	\qwhereq
	x_m = \dotp{f}{\psi_m}.
}
This count is also called the $\lzero$ pseudonorm:
\eq{
	\normz{x} \eqdef \Jz(f).
}
It measures sparsity directly, without accounting for the magnitudes of the nonzero coefficients.

Natural images usually have many nonzero coefficients, but can often be approximated by a function $f_M$ with a small count $M=\Jz(f_M)$. Retaining coefficients above a threshold gives a best $M$-term approximation, with $M$ determined by that threshold:
\eq{
	f_M = \sum_{ |\dotp{f}{\psi_m}|>T } \dotp{f}{\psi_m} \psi_m
	\qwhereq
	M = \#\enscond{m}{|\dotp{f}{\psi_m}|>T}.
}
For suitable wavelet bases and the bounded-variation image classes specified in Section \ref{sec-signal-models}, the error $\norm{f-f_M}$ satisfies quantitative decay estimates. These estimates motivate approximating natural images $f$ by functions with small $\Jz$.

Figure \ref{fig-sparsity-wav} displays a natural image in the wavelet basis $\psi_m = \psi_{j,n}^\om$, indexed by $m=(j,n,\om)$. Most coefficients $\dotp{f}{\psi_m}$ have small magnitude, so retaining the few large coefficients captures much of the image.

\myfigure{
\tabdeux{
\image{variational}{.3}{sparsity-image}&
\image{variational}{.3}{sparsity-wav}\\
Image $f$ & Coefficients $\dotp{f}{\psi_m}$
}
}{ 
    Most wavelet coefficients of a natural image have small magnitude.
}{fig-sparsity-wav}


                                                                      
\subsection{Convex Relaxation}

The ideal sparsity prior $\Jz$ is nonconvex, which makes optimization difficult. For example, if $f$ and $g$ have disjoint nonempty coefficient supports in $\Bb$, then $\Jz( (f+g)/2 ) = \Jz(f)+\Jz(g)$, violating convexity of $\Jz$.

In a general inverse problem, minimizing $\Jz$ entails a combinatorial search over possible coefficient supports.

To approximate the ideal prior $\Jz$, consider the $\ell^q$ family with $q>0$:
\eq{
	J_q(f) = \displaystyle \sum_m |\dotp{f}{\psi_m}|^q.
}
As shown in Figure \ref{fig-sparsity-lq}, the unit balls in $\RR^2$ become concentrated near the coordinate axes as $q\downarrow0$. For each fixed finite-dimensional vector, $J_q(f)\to\Jz(f)$; the limiting shape of these bounded unit balls should not be confused with the unbounded set $\{x:\norm{x}_0\leq1\}$. Small values of $q$ favor sparsity.


\myfigure{
\image{variational}{1}{sparsity-lq-balls}
}{ 
	$\ell^q$ balls $\enscond{x}{J_q(x) \leq 1}$ for varying $q$.
}{fig-sparsity-lq}


The prior $J_q$ is convex exactly when $q \geq 1$. The smallest exponent yielding a convex prior, $q=1$, gives the $\lun$ prior $\Ju$:
\eql{\label{eq-dfn-lun-prior}
	\Ju(f) = \norm{(\dotp{f}{\psi_m})}_1 = \sum_m |\dotp{f}{\psi_m}|.
}
In the following, we consider discrete orthonormal bases $\Bb = \{\psi_m\}_{m=0}^{N-1}$ of $\RR^N$.


                                                                      
\subsection{Sparse Regularization and Thresholding}

Given an orthonormal basis $\{\psi_m\}_m$ of $\RR^N$, regularization-based denoising \eqref{eq-regul-denoising} can be written using the sparsity priors $\Jz$ and $\Ju$ as
\eq{
	f^\star \in \uargmin{g \in \RR^N} \frac{1}{2} \norm{f-g}^2 + \la J_q(g)
}
for $q=0$ or $q=1$, writing $J_0=\Jz$. Orthonormality separates the objective into coefficientwise terms:
\eq{
	f^\star = \sum_m x^\star_m \psi_m
}
\eq{
	\qwhereq
	x^{\star} \in \uargmin{y \in \RR^N} \sum_m \frac{1}{2} |x_m-y_m|^2 + \la |y_m|^q
}
Here $x_m \eqdef \dotp{f}{\psi_m}$ and $y_m \eqdef \dotp{g}{\psi_m}$. For $q=0$, we adopt the convention
\eq{
	\foralls u \in \RR, \quad
	|u|^0=
	\choice{
		0 \qifq u=0, \\ 1 \quad \text{otherwise.} 
	}
}
Each coefficient of the denoised image solves a one-dimensional optimization problem
\eql{\label{eq-var-1D-spars}
	x_m^\star \in \uargmin{u \in \RR}  \frac{1}{2} |x_m-u|^2 + \la |u|^q
}
The following proposition gives a closed-form solution by thresholding.

\begin{prop}\label{prop-equiv-sparse-thresh}
One may choose the minimizer
\eql{\label{eq-equiv-sparse-thresh}
	x^{\star}_m = S_T^q( x_m )
	\qwhereq 
	\choice{
		T = \sqrt{2 \la} \qforq q=0,\\
		T = \la \qforq q=1,
	}
}
where 
\eql{\label{eq-hard-thresh}
	\foralls u \in \RR, \quad
	S_T^0(u) \eqdef
	\choice{
		0 \qifq |u| < T, \\
		u \quad\text{otherwise}	
	}
} 
is hard thresholding, with the alternative tie convention discussed below, and
\eql{\label{eq-soft-thresh}
	\foralls u \in \RR, \quad
	S_T^1(u) \eqdef  \sign(u) ( |u|-T )_+
} 
is the soft thresholding introduced